%Chargement des paquets
\usepackage{amsmath}
\usepackage{amsfonts}
\usepackage{amssymb}
\usepackage{enumerate}
\usepackage{mathtools}
\usepackage{bbm}
\usepackage{xparse, etoolbox}
\usepackage{enumerate}
\usepackage{enumitem}
\usepackage{mathabx}
\usepackage{babel}[french]

%environment
\usepackage{amsthm}
\usepackage{keytheorems}
\usepackage{hyperref} 

%Interface théorème
\newkeytheorem{lem}[
	name = Lemme
]

\newkeytheorem{prop}[
	name = Proposition
]

\newkeytheorem{thm}[
	name = Théorème
]

\newkeytheorem{coro}[
	name = Corollaire
]

\renewcommand*{\proofname}{Démonstration}

\theoremstyle{definition}
\newtheorem{defn}{Définition}

\theoremstyle{definition}
\newtheorem*{nota}{Notation}

\theoremstyle{definition}
\newtheorem*{limi}{Liminaire}

\theoremstyle{remark}
\newtheorem{rmq}{Remarque}

\theoremstyle{plain}
\newtheorem*{fait}{Fait}


%Ensembles de nombres
\newcommand{\N} {\mathbb{N}}
\newcommand{\Ne}{\N^\ast}
\newcommand{\Z} {\mathbb{Z}}
\newcommand{\Q} {\mathbb{Q}}
\newcommand{\R} {\mathbb{R}}
\newcommand{\Rb}{\overline{\mathbb{R}}}
\newcommand{\Rp}{\R_+}
\newcommand{\Rm}{\R_-}
\newcommand{\K} {\mathbb{K}}
\newcommand{\Ket} {\mathbb{K^{\ast}}}
\newcommand{\Cx}{\mathbb{C}}

%Ensembles, tribus
\newcommand{\A}{\mathbb{A}}
\renewcommand{\P}{\mathcal{P}}
\newcommand{\T}{\mathcal{T}}
\newcommand{\F}{\mathcal{F}}
\newcommand{\C} {\mathcal{C}}
\newcommand{\ouv}{\mathcal{O}}
\newcommand{\B}{\mathcal{B}}
\newcommand{\M}{\mathcal{M}}
\newcommand{\etag}{\mathcal{E}}
\newcommand{\etagOp}{\etag^{+}(\Omega)}
\newcommand{\etagO}{\etag(\Omega)}
%%Lp avant quotientage
\newcommand{\Lic}{\mathcal{L}^1}
\newcommand{\Lpc}{\mathcal{L}^p}
\newcommand{\Linfc}{\mathcal{L}^{\infty}}
\newcommand{\Lifull}{\Lic(\Omega, \B, m)}
\newcommand{\Lpfull}{\Lpc(\Omega, \B, m)}
\newcommand{\Linffull}{\Linfc(\Omega, \B, m)}
%%Lp après quotientage
\newcommand{\Li}{L^1}
\newcommand{\Ld}{L^2}
\newcommand{\Lp}{L^p}
\newcommand{\Linf}{L^{\infty}}
\newcommand{\Listd}{\Li(\Omega, m)} 
\newcommand{\Ldstd}{\Ld(\Omega, m)} 
\newcommand{\Lpstd}{\Lp(\Omega, m)} 

%Quelques opérateurs
\DeclareMathOperator{\codim}{codim}
\DeclareMathOperator{\rg}{rg}
\DeclareMathOperator{\tr}{tr}
\DeclareMathOperator{\vect}{Vect}
\DeclareMathOperator{\ima}{Im}
\DeclareMathOperator{\id}{Id}
\DeclareMathOperator{\sign}{sign}
\DeclareMathOperator{\ext}{ext}
\newcommand{\ind}{\mathbbm{1}}
\newcommand*{\spr}[2]{\left(\,#1\,\middle|\,#2\,\right)}
\newcommand{\interior}[1]{%
  {\kern0pt#1}^{\mathrm{o}}%
}
\DeclareMathOperator{\card}{card}
\DeclareMathOperator{\ppcm}{ppcm}

%Commandes de pure flemme
\newcommand{\veps}{\varepsilon}
\newcommand{\intab}{\int_{a}^{b}}
\newcommand{\intR}{\int_{-\infty}^{\infty}}
\newcommand{\intinf}{\int_{- \infty}^{+ \infty}}
\newcommand{\equi}{\Leftrightarrow}
\newcommand{\Kev}{$\K$-espace vectoriel }
\newcommand{\Kevs}{$\K$-espaces vectoriels }
\newcommand{\Rev}{$\R$-espace vectoriel }
\newcommand{\Revs}{$\R$-espaces vectoriels }
\newcommand{\Cev}{$\Cx$-espace vectoriel }
\newcommand{\Cevs}{$\Cx$-espaces vectoriels }
\newcommand{\evn}{(E, \norm{.})}
\newcommand{\interf}[2]{[#1,#2]}
\newcommand{\limb}{\lim\limits_}
\newcommand{\lebn}{\lambda_n}
\newcommand{\pp}{\text{~presque partout}}
\newcommand{\derivp}[2]{\frac{\partial #1}{\partial #2}}
\newcommand{\famiu}{(u_i)_{i \in I}}
\newcommand{\sev}{\text{~s.e.v.}}
\newcommand{\Mnp}{M_{n,p}(\K)}
\newcommand{\Mn}{M_{n}(\K)}
\newcommand{\tq}{\text{~t.q.~}}
\newcommand{\mq}{~montrer~que~}
\newcommand{\et}{\quad \text{et} \quad}
\newcommand{\ou}{\text{~ou~}}
\newcommand{\Kx}{\K[X]}


%Commandes de gars chelou
\renewcommand{\subset}{\subseteq}

%Commande restriction, ty duckduckgo
\def\restr#1#2{\mathchoice
              {\setbox1\hbox{${\displaystyle #1}_{\scriptstyle #2}$}
              \restraux{#1}{#2}}
              {\setbox1\hbox{${\textstyle #1}_{\scriptstyle #2}$}
              \restraux{#1}{#2}}
              {\setbox1\hbox{${\scriptstyle #1}_{\scriptscriptstyle #2}$}
              \restraux{#1}{#2}}
              {\setbox1\hbox{${\scriptscriptstyle #1}_{\scriptscriptstyle #2}$}
              \restraux{#1}{#2}}}
\def\restraux#1#2{{#1\,\smash{\vrule height .8\ht1 depth .85\dp1}}_{\,#2}} 

%Quelques normes
\DeclarePairedDelimiter\abs{\lvert}{\rvert}
\DeclarePairedDelimiter\labs{\bigl\lvert}{\bigr\rvert}
\DeclarePairedDelimiter\norm{\lVert}{\rVert}
\DeclarePairedDelimiterXPP\normi[1]{}\lVert\rVert{_1}{\ifblank{#1}{\:\cdot\:}{#1}}
\DeclarePairedDelimiterXPP\normd[1]{}\lVert\rVert{_2}{\ifblank{#1}{\:\cdot\:}{#1}}
\DeclarePairedDelimiterXPP\normp[1]{}\lVert\rVert{_p}{\ifblank{#1}{\:\cdot\:}{#1}}
\DeclarePairedDelimiterXPP\normq[1]{}\lVert\rVert{_q}{\ifblank{#1}{\:\cdot\:}{#1}}
\DeclarePairedDelimiterXPP\norminf[1]{}\lVert\rVert{_\infty}{\ifblank{#1}{\:\cdot\:}{#1}}

%Bracket en angle
\DeclarePairedDelimiter\angl{\langle}{\rangle}

%Set, parenthèses
\newcommand{\set}[1]{\{ #1 \}}
\newcommand{\bigset}[1]{\bigl\{ #1 \bigr\}}
\newcommand{\Bigset}[1]{\Bigl\{ #1 \Bigr\}}
\newcommand{\bigpar}[1]{\bigl( #1 \bigr)}
\newcommand{\Bigpar}[1]{\Bigl( #1 \Bigr)}

%Opitmisation des opérateurs de comparaison
\DeclarePairedDelimiter\tio{o (}{)}
\DeclarePairedDelimiter\gro{O (}{)}

\renewcommand{\L}{\mathcal{L}}

\pagestyle{fancy}

\fancyhead[C]{Théorème d'inversion locale}
\fancyhead[R]{}

\begin{document}

\underline{Source :} Analyse - Gourdon - page 341 \newline

\underline{Leçons :}
\begin{itemize}
\item 215 : Applications différentiables définies sur un ouvert de Rn. Exemples et applications.
\item 214 : Théorème d’inversion locale, théorème des fonctions implicites. Illustrations en analyse et en géométrie.
\end{itemize}
\vspace{5pt}

\underline{Introduction :} ce théorème repose sur une idée simple, on considère une fonction $f$ entre deux espaces vectoriels et on
suppose que sa différentielle en un point est inversible. En d'autres termes, $f$ peut être localement approximée par un isomorphisme
d'espaces vectoriels. On peut donc légitimement supposer que, au moins localement, $f$ est inversible.
C'est ce que l'on démontre en imposant la continuité de la différentielle (ce qui permet de contrôler l'aspect local du résultat). 
La bicontinuité de la différentielle permet d'obtenir une inverse locale continue.

\begin{defn}
Soient $E$ et $F$ deux \Rev, $U$ un ouvert de $E$ et $a \in U$. Une application $f:U \to F$ est dite différentiable au point $a$, s'il existe
une application linéaire continue de $E$ dans $F, L \in \L_C(E,F)$, et telle que :
\[ f(a+h) = f(a) + L(h) + o(\norm{h}) \text{ lorsque } h \to 0 . \]
\end{defn}

\begin{lem}
Si une telle application linéaire existe, elle est unique. On l'appelle différentielle de $f$ en $a$ et on la note $df_a$.
\end{lem}

\begin{proof}
Prenons deux telles applications linéaires notées $L$ et $L'$. On pose $u = L - L'$ et on veut montrer que $u=0$.
On prend $s$ un vecteur de norme 1 dans $E$ et $V$ un voisinage de 0 dans $F$, on a $u \in o(\norm{h})$, donc on peut trouver
$\alpha > 0$ tel que si $\norm{h} \leq \alpha$ alors $\dfrac{u(h)}{\norm{h}} \in V$. Or, par linéarité de $u$ :
\[ u(s) = \dfrac{u(\alpha s)}{\alpha} = \dfrac{u(\alpha s)}{\norm{\alpha s}} \in V. \]
Ainsi, $u(s)$ est dans tous les voisinages de $0$, donc $u(s) = 0$ par séparation de $F$. Autrement dit, $u$ est une application linéaire nulle
sur toute la boule unité de $E$, elle est donc nulle sur tout $E$.
\end{proof}

\begin{nota}
Dans la suite, pour $r>0$ on note $B_r$ la boule ouverte de centre 0 et de rayon $r$ et $\widebar{B_r}$ la boule fermée correspondante.
$\L_C(E,F)$ est l'ensemble des applications linéaires continues de $E$ dans $F$. On norme ce \Rev par 
\[ \normt{u} = \sup_{x \in B_1} \norm{u(x)} . \]
\end{nota}

\begin{defn}
Si $f$ est différentiable en tout point de $U$, on dit que $f$ est différentiable sur $U$ et l'application :
\[ \begin{array}{c c c c}
df : & U & \to & \L_C(E,F) \\
& a & \mapsto & df_a
\end{array} \]
est appelée application différentielle de $f$. Si $df$ est continue, on dit que $f$ est de classe $\C^1$ sur $U$.
\end{defn}

\begin{lem}
\label{lem:inv}
Soient $E$ un espace de Banach et $u \in \L_C(E)$ telle que $\normt{u} <1$. Alors $\id - u$ est inversible et son inverse est :
\[ \sum_{n=0}^{\infty} u^n \in \L_C(E) . \]
\end{lem}

\begin{proof}
Comme $\L_C(E)$ est une algèbre normée, on a $\forall n \in \N, \normt{u^n} \leq \normt{u}^n < 1$ donc la série $\sum u_n$ converge
absolument. On vérifie ensuite le résultat par calcul :
\[ (\id-u) \left ( \sum_{n=0}^{\infty} u^n \right ) = \sum_{n=0}^{\infty} u^n - \sum_{n=1}^{\infty} u^n = \id, \]
et de même pour $\left ( \sum_{n=0}^{\infty} u^n \right ) (\id - u)$.
\end{proof}

\begin{thm}[Inversion locale]
Soient $E$ et $F$ deux espaces de Banach, $U$ un ouvert de $E$ et $f:U \to F$ une application de classe $\C^1$.
On suppose qu'il existe $a \in U$ tel que $df_a$ soit un isomorphisme bicontinu de $E$ sur $F$. 
Alors il existe un voisinage ouvert $V$ de $a$ et un voisinage ouvert $W$ de $f(a)$ tels que :
\begin{enumerate}[label=(\roman*)]
\item la restriction $\restr{f}{V}$ de $f$ à $V$ est une bijection de $V$ sur $W$ ;
\item l'application inverse $g:W \to V$ est continue ;
\item $g$ est de classe $\C^1$ et, pour tout $x \in V, dg_{f(x)} = df_x^{-1}$.
\end{enumerate}
\end{thm}

\begin{proof}
On commence par simplifier le problème, on peut en effet toujours se ramener au cas $a=0, f(a)=0$ et $df_0 = \id$ (et donc $E=F$)
en considérant l'application $x \mapsto df_a^{-1}(f(a+x) -f(a))$ définie sur $E$. \\
Comme $f$ est de classe $\C^1$ sur $U$, on peut trouver $r>0$ tel que $\widebar{B_r} \subset U$ et, pour tout $x \in \widebar{B_r}$ :
\[ \normt{df_x - df_0} = \normt{df_x - \id} \leq 1/2 . \]
Pour $x \in \widebar{B_r}$, on peut écrire $df_x = \id - (\id - df_x)$ et comme $\normt{\id - df_x} < 1, df_x$ est inversible d'inverse 
$df_x^{-1} = \sum_{n=0}^{\infty} (\id - df_x)^n$ (voir lemme $\ref{lem:inv}$) et de plus :
\[ \normt{df_x} \leq \sum_{n=0}^{\infty} \normt{u}^n \leq \sum_{n=0}^{\infty} \dfrac12^n = 2 . \qquad (*) \]
Autrement dit $df_x$ est un isomorphisme bicontinu.

\begin{enumerate}[label=(\roman*)]

\item 

On veut montrer que $f$ admet localement une fonction inverse, pour cela on va démontrer que :
\[ \forall y \in B_{r/2}, \exists ! x \in B_r \tq y=f(x). \]
On fixe donc $y \in B_{r/2}$ et on considère la fonction $h : B_r \to E, x \mapsto y + x - f(x)$. 
Le problème que l'on se pose revient à démontrer que $h$ admet un unique point fixe. 
$h$ est de classe $\C^1$, et, pour tout $x \in B_r, \normt{dh_x} = \normt{\id - df_x} \leq 1/2$, 
donc, d'après l'inégalité des accroissements finis :
\[ \forall x, x' \in \widebar{B_r}, \norm{h(x)-h(x')} \leq \dfrac12 \norm{x-x'} . \qquad (**) \]
En particulier, pour tout $x \in \widebar{B_r}, \norm{x-f(x)} = \norm{h(x) - h(0)} \leq \norm{x}/2$, donc :
\[ \forall x \in \widebar{B_r}, \norm{h(x)} \leq \norm{y} + \norm{x-f(x)} \leq \norm{y} + \dfrac12 \norm{x} 
	\leq \dfrac{r}2 +\dfrac{r}2 \leq r . \]
Ainsi, $h$ est une fonction de $\widebar{B_r}$ dans $B_r \subset \widebar{B_r}$ et elle est $1/2$-contractante. 
Donc, selon le théorème du point fixe de Banach-Picard, elle admet un unique point fixe $x \in B_r$. \\
En résumé, nous avons montré que pour tout $y \in B_{r/2}$ il existe un unique $x \in B_r$ tel que $f(x)=y$. 
Ainsi, $V = f^{-1}(B_{r/2}) \cap B_r$ est un ouvert qui contient $0$ (car $f$ est continue et $f(0)=0$). 
En posant $W = B_{r/2}$, on a bien démontré que $\restr{f}{V}$ est une bijection de $V$ sur $W$.

\item

Notons $g: W \to V$ l'application inverse de $\restr{f}{V}$. On considère l'application $h$ du point précédent pour $y=0$, soit 
$h : B_r \to E, x \mapsto x - f(x)$. On remarque que $x = f(x) + h(x)$ pour tout $x$ dans $B_r$. Ainsi :
\[ \forall x, x' \in B_r, \norm{x-x'} \leq \norm{h(x)-h(x')} + \norm{f(x)-f(x')} \leq \dfrac12 \norm{x-x'} + \norm{f(x)-f(x')}, \]
donc $\norm{x-x'} \leq 2\norm{f(x)-f(x')}$. On en déduit :
\[ \forall y, y' \in W, \norm{g(y) - g(y')} \leq 2 \norm{f(g(y)) - f(g(y'))} = 2 \norm{y-y'} \qquad (***) \]
autrement dit, $g$ est lipschitzienne, donc continue.

\item

On fixe $x \in V$ et on pose $y = f(x) \in W$. Soit $w$ tel que $y+w \in W$, et on note $v = g(y+w) - g(y)$.
On a $\norm{v} \leq 2 \norm{w}$ d'après (***), et :
\begin{align*}
\delta(w) & = g(y+w) - g(y) - df_x^{-1}(w) \\
& = v - df_x^{-1}(f(x+v)-f(x)) \\
& = - df_x^{-1}(f(x+v)-f(x)-df_x(v)) 
\end{align*}
Or $\normt{df_x^{-1}} \leq 2$ d'après (*) donc :
\[ \norm{\delta(w)} \leq 2 \norm{f(x+v)-f(x)-df_x(v)} = 2 \norm{v} \veps(v) , \]
avec $\lim_{v \to 0} \veps(v) = 0$, et :
\[ \norm{\delta(w)} \leq 4 \norm{w} \veps(g(y+w)-g(y)) = 4 \norm{w} \veps'(w) . \]
Par continuité de la fonction $g$ en 0, on a $\lim_{w \to 0} \veps'(w) = 0$, donc $\norm{\delta(w)} = o(\norm{w})$.
La fonction $g$ est donc différentiable en $y$ et $dg_y = df_x^{-1}$. \\
Enfin, comme $f$ est de classe $\C^1$, que la fonction $g$ est continue et que la fonction $u \mapsto u^{-1}$ (définie sur $\L_C(E)$)
est continue, l'application $y \mapsto dg_y = df_{g(y)}^{-1}$ est continue, donc $g$ est de classe $\C^1$.

\end{enumerate}

\end{proof}

\begin{defn}
Soient $E_1, \dots, E_n, F$ des e.v.n. et une application $f: \Omega \subset E_1 \times \dots \times E_n \to F$, où $\Omega$ 
est un ouvert de $E_1 \times \dots \times E_n $. Soit $a=(a_1, \dots, a_n) \in \Omega$. 
Pour $i \in \set{1, \dots, n}$, la fonction $f_i : x \mapsto f(a_1, \dots, a_{i-1},x,a_{i+1}, \dots, a_n)$ est définie sur un voisinage de $a_i$ 
dans $E_i$. Si elle est différentiable en $a_i$, on dit que $f$ admet une différentielle partielle d'indice $i$, et on appelle différentielle
partielle d'indice $i$ la différentielle $df_{i,a_i}$, notée $\partial_i f_{(a_1, \dots, a_n)}$.
\end{defn}

\begin{rmq}
Si une telle fonction $f$ est différentiable en $a=(a_1, \dots, a_n)$ alors $\partial_i f_{(a_1, \dots, a_n)}$ existe pour tout $i$ et :
\[ \forall h \in E_i, \quad \partial_i f_{(a_1, \dots, a_n)}(h) = df_a(0, \dots, 0, h, 0, \dots, 0) , \]
ce qui entraîne :
\[ \forall h=(h_1, \dots, h_n) \in E_1 \times \dots \times E_n, \quad df_a(h) = \sum_{i=1}^n \partial_i f_{(a_1, \dots, a_n)}(h_i) . \]
\end{rmq}

\begin{thm}[Fonctions implicites]
Soient $E, F, G$ trois espaces de Banach, $U$ un ouvert de $E$ et $V$ un ouvert de $F$.
On considère l'application :
\[ \begin{array}{c c c c}
f : & U \times V & \to & G \\
& (x,y) & \mapsto & f(x,y)
\end{array} \]
de classe $\C^1$. On suppose de plus que pour tout $(a,b) \in U \times V, \partial_2 f_{(a,b)}$ est un isomorphisme bicontinu de $F$ sur
$G$. Alors il existe :
\begin{enumerate}[label=(\roman*)]
\item un voisinage ouvert $U'$ de $a$ avec $U' \subset U$,
\item un voisinage ouvert $W$ de $f(a,b)$,
\item un voisinage ouvert $\Omega$ de $(a,b)$ avec $\Omega \subset U \times V$,
\item une fonction $\varphi : U' \times W \to V$ de classe $\C^1$ et vérifiant,
\item pour tout $x \in U'$, pour tout $z \in W$, il existe une unique solution $y$ de $f(x,y)=z$, 
avec $(x,y) \in \Omega$ et on a $y=\varphi(x,z)$.
\end{enumerate}
En particulier, $f(x, \varphi(x,z))=z$ pour tout $(x,z) \in U' \times W$.
\end{thm}

\begin{proof}
On introduit la fonction :
\[ \begin{array}{c c c c}
F : & E \times F & \to E \times G \\
& (x,y) & \mapsto & (x, f(x,y))
\end{array} \]
Cette fonction est de classe $\C^1$ et :
\[ dF_{(a,b)}(h_1,h_2) = (h_1, df_{(a,b)}(h_1, h_2)) = (h_1, \partial_1 f_{(a,b)}(h_1) + \partial_2 f_{(a,b)}(h_2)) . \]
Ainsi, $dF_{(a,b)}$ est un isomorphisme bicontinu de $ E \times F$ sur $E \times G$ et son isomorphisme inverse est donné par :
\[ (l_1, l_2) \mapsto (l_1, \partial_2 f_{(a,b)}^{-1}(l_2 - \partial_1 f_{(a,b)}(l_1)) . \]
On peut donc appliquer le théorème d'inversion locale, ce qui entraîne l'existence d'un voisinage ouvert $\Omega$ de $(a,b)$
et d'un voisinage ouvert $\Gamma$ de $(a,f(a,b))$ tels que $F$ soit un $\C^1$-difféomorphisme de $\Omega$ sur $\Gamma$.
Quitte à restreindre $\Gamma$  on peut même écrire $\Gamma = U' \times W$ où $U'$ est un voisinage ouvert de $a$ et $W$ un 
voisinage ouvert de $f(a,b)$. \\
Comme $F(x,y)=(x, f(x,y))$, le difféomoprhisme inverse $F^{-1} : \Gamma \to \Omega$ s'écrit sous la forme :
\[ F^{-1}(x,z) = (x, \varphi(x,z)) , \]
avec $\varphi$ de classe $\C^1$. On en déduit que, pour tout $(x,z) \in \Gamma = U' \times W$, il existe un unique $y$ tel que
$(x,y) \in \Omega$ et $f(x,y)=z$, et de plus $y=\varphi(x,z)$. 
\end{proof}

\newpage

\section*{Annexe : inégalités des accroissements finis}

\begin{thm}[cas $\R \to F$]
Soient deux réels $a < b$ et deux fonctions $f : [a,b] \to F$ et $g : [a,b] \to \R$, toutes deux continues sur $[a,b]$ et dérivables sur
$]a,b[$. Si :
\[ \forall t \in ]a,b[, \quad \norm{f'(t)} \leq g'(t) \]
alors :
\[ \norm{f(b)-f(a)}\leq g(b) - g(a) . \]
\end{thm}

% preuve 1, vraiment pas fan (source wikiversité)

%\begin{proof}
%Pour $\veps > 0$ et $u \in ]a,b[$, on définit sur $[u,b]$ la fonction $h$ par :
%\[ h(t) = \norm{f(t)-f(a)} - g(t) - \veps t . \]
%Comme $h$ est continue, on note $c$ le point de $[u,b]$ pour lequel $h$ atteint son minimum. On va démontrer que $c=b$. \\
%Sinon, l'intervalle $]c,b[$ est non vide, on peut donc trouver $t \in ]c,b[$ tel que :
%\[ \norm*{\dfrac{f(t)-f(c)}{t-c}} - \dfrac{\veps}{2} < \norm{f'(c)} \leq g'(c)  < \dfrac{g(t)-g(c)}{t-c} > + \dfrac{\veps}{2} \]
%soit :
%\[ \norm{f(t)-f(c)} < g(t) - g(c) + \veps(t-c), \]
%et donc :
%\begin{align*} 
%h(t) - h(c) & = \norm{f(t) - f(a)} - g(t) - \veps t - ( \norm{f(c) - f(a)} - g(c) - \veps c ) \\
%& = \norm{f(t) - f(c) + f(c) - f(a)} - \norm{f(c) - f(a)} - g(t) + g(c) - \veps (t-c) \\
%& \leq \norm{f(t) - f(c)} + \norm{f(c) - f(a)} - \norm{f(c) - f(a)} - g(t) + g(c) - \veps (t-c) \\
%& < 0
%\end{align*}
%soit $h(t) < h(c)$, ce qui contredit le choix de $c$. Donc $c=b$, en particulier $h(b) \leq h(u)$ i.e. :
%\[ \norm{f(b) - f(a)} \leq \norm{f(u) - f(a)} + g(b) - g(u) + \veps(b-u) . \]
%On conclut en passant à la limite pour $\veps \to 0^{+}$ et $u \to a^{+}$.
%\end{proof}

% preuve 2 du Gourdon, plus longue et même idée mais ça fait moins parachuté

\begin{proof}
On suppose d'abord que $norm{f'(t)} < g'(t)$ sur $]a,b[$. Donc, pour tout $x \in ]a,b[$, on a :
\[ \lim_{\substack{t \to x \\ t > x}} \norm*{\dfrac{f(t)-f(x)}{t-x}} - \dfrac{g(t)-g(x)}{t-x} < 0 , \]
on peut donc trouver un intervalle de borne inférieure $x$ sur laquelle cette inégalité est vraie, soit :
\[ \forall x \in ]a,b[, \exists y > x, \forall t \in ]x, y], \norm*{\dfrac{f(t)-f(x)}{t-x}} < \dfrac{g(t)-g(x)}{t-x} , \]
avec égalité des numérateurs pour $t=x$ soit :
\[ \forall x \in ]a,b[, \exists y > x, \forall t \in [x, y], \norm*{f(t)-f(x)} \leq g(t)-g(x) . \qquad (*) \]
Considérons maintenant un compact $[c,d] \subset ]a,b[$ et montrons que :
\[ \norm{f(d)-f(c)} \leq g(d)-g(c) . \qquad (**) \]
Pour cela, on considère l'ensemble :
\[ \Delta = \set{ z \in ]c,d], \forall t \in [c, z], \norm{f(t)-f(c)} \leq g(t) - g(c) } . \]
$\Delta$ est non vide d'après (*), il admet donc une borne supérieure que l'on note $m$. 
On veut montrer que $m=d$ ce qui prouvera que l'inégalité (**) est vraie.
Comme $f$ et $g$ sont continues, on a $\norm{f(m)-f(c)} \leq g(m) - g(c)$. 
Supposons que $m < d$, alors d'après (*) :
\[ \exists \gamma \in ]m, d], \forall t \in [m, \gamma], \norm{f(t)-f(m)} \leq g(t) - g(m) , \]
donc 
\[ \forall t \in [m, \gamma], \norm{f(t)-f(c)} \leq \norm{f(t)-f(m)} + \norm{f(m)-f(c)} \leq g(t) - g(c) . \]
Autrement dit, $\gamma \in \Delta$ ce qui est absurde car $\gamma > m = \sup{\Delta}$. \\
On a donc prouvé l'inégalité (**), ce qui permet de déduire, en faisant tendre $c \to a$, $d \to b$, et
en vertue de la continuité des fonctions $f$ et $g$ :
\[ \norm{f(b)-f(a)}\leq g(b) - g(a) . \]

On se ramène au cas général en définissant, pour $\veps >0$, la fonction $g_{\veps}$ par :
\[ g_{\veps}(t) = g(t) + \veps t . \]
On a bien $\norm{f'(t)} < g_{\veps}'(t)$ sur $]a,b[$ donc, d'après ce que l'on vient de démontrer :
\[ \norm{f(b)-f(a)}\leq g(b) - g(a) + \veps (b-a) . \]
On conclut en faisant tendre $\veps \to 0$.
\end{proof}

\begin{thm}[cas $E \to F$]
Soit $f:U \to F$, où $E$ et $F$ sont deux e.v.n et $U$ un ouvert de $E$. 
Soient $a,b \in U$ tels que le segment $[a,b]$ soit inclus dans $U$. 
Si $f$ est continue sur $[a,b]$, différentiable en tout point de $]a,b[$ et s'il existe $M>0$ tel que $\normt{df_c} \leq M$ pour tout 
$c \in ]a,b[$, alors :
\[ \norm{f(b)-f(a)} \leq M \norm{b-a} . \]
\end{thm}

\begin{proof}
La fonction $g : [0,1] \to F, t \mapsto f(a+t(b-a))$ est continue sur [0,1], dérivable sur ]0,1[ et :
\[ \forall t \in ]0,1[, \quad g'(t) = df_{a+t(b-a)}(b-a) , \]
autrement dit, $\norm{g'(t)} \leq M \norm{b-a}$ sur $]0,1[$. Donc, en appliquant le théorème dans le cas $\R \to F$, on conclut que :
\[ \norm{g(1) - g(0)} = \norm{f(b)-f(a)} \leq M \norm{b-a} . \] 
\end{proof}

\begin{rmq}
On peut étendre ce résultat à toute partie convexe de $U$, voir à tout étoilé de $U$ en fixant un point de référence.
\end{rmq}

\end{document}